\begin{center}
Problema D

Balança do Festival

{\small Nome do arquivo: D.c, D.cpp, D.java, ou D.py}

{\large Timelimit: $1$}

{\tiny Por Athos José}


\end{center}

Na cidade das Abóboras, Zé do Grafo é famoso por fazer suas peripécias com grafo, mas agora ele quer cultivar as mais saborosas abóboras da região. Com a chegada do \textit{Concurso das Abóboras}, ela precisa escolher algumas de suas abóboras para exibir no festival. No entanto, existe uma regra: a soma dos pesos das abóboras escolhidas deve ser um valor $P$ estabelecido pelo festival.

Zé do Grafo já colheu algumas abóboras e está determinado a participar, mas precisa da sua ajuda para decidir se ele conseguirá escolher algumas abóboras que atinjam exatamente o peso estabelecido, sem passar nem faltar. Será que você consegue ajudá-lo dizendo se ele tem essas abóboras?

\vspace{0.3cm}
\textbf{Entrada}

A primeira linha contém dois inteiros, $N$, $P$, representando respectivamente o número de abóboras $(1 \leq N \leq 10^4)$ e o peso para competição $(1 \leq P \leq 10^4)$. A segunda linha contém $N$ inteiros $a_1, \dots, a_N$ $(1 \leq a_i \leq 100)$, onde $a_i$ indica o peso da abóbora.

\vspace{0.3cm}
\textbf{Saída}

Imprima ``SIM'' se for possível escolher algumas abóboras de modo que a soma dos pesos delas seja exatamente $P$. Caso contrário, imprima ``NAO''.

\begin{table}[h]
    \centering
    \begin{tabular}{|p{7cm}|p{7cm}|}
       \hline
        \begin{center}
            \vspace{-0.3cm}
            \textbf{Exemplos de Entrada}
            \vspace{-0.3cm}
        \end{center} 
         &  
        \begin{center}
            \vspace{-0.3cm}
            \textbf{Exemplos de Saída}
            \vspace{-0.3cm}
        \end{center} \\ \hline
         \texttt{5 9} & \texttt{SIM} \\  
         \texttt{3 1 4 2 5} & \\
         \hline
         \texttt{4 11} & \texttt{NAO} \\  
         \texttt{8 4 6 2} & \\
         \hline
    \end{tabular}
    \caption{Exemplos de entradas e saídas}
\end{table}

\newpage

\begin{center}
\noindent
\lstinputlisting{abobora_24/D/D4.cpp}
\end{center}